\documentclass[10pt]{scrartcl}
\usepackage[mdthm,colorsec,noauthor,noasy]{evan}
\usepackage{handout}
\newcommand{\ip}[2]{\langle #1,#2\rangle}
\DeclareMathOperator{\tr}{tr}\DeclareMathOperator{\nul}{null}\DeclareMathOperator{\range}{range}\DeclareMathOperator{\spn}{span}
\newcommand{\LL}{\mathcal L}\newcommand{\PP}{\mathcal P}\newcommand{\MM}{\mathcal M}
\title{Math 110 --- Week 9}
\subtitle{Trace and inner products (Notes Ch.~16--17, LADR 8D, 6A, HW9)}
\author{}\date{}\ihead{}
\begin{document}
\maketitle
\section{Facts used from earlier weeks}
$V$ is finite-dimensional over $\FF$ ($\RR$ or $\CC$), $n=\dim V$. Shorthand: ``$S$ kills $v$'' means $Sv=0$, and ``$q$ kills $T$'' means $q(T)=0$. From Week 8: over $\CC$, $V=\bigoplus_iG(\lambda_i,T)$ with $d_i=\dim G(\lambda_i,T)$ (Prop.~73), $\chi_T=\prod_i(x-\lambda_i)^{d_i}$ (Prop.~74), and $T$ has a basis in which its matrix is upper-triangular with each $\lambda_i$ appearing $d_i$ times on the diagonal (LADR 8.37; a Jordan basis, Prop.~76, is one such basis).

\begin{enumerate}[label=(T\arabic*),leftmargin=2.6em,itemsep=8pt]
  \ii \tool{Change of basis.} If column $j$ of $S$ holds the $B$-coordinates of the $j$-th vector of $B'$, then $\MM_{B'}(T)=S^{-1}\MM_B(T)S$. So similar matrices are exactly the matrices of one operator in two bases.
  \why{$S$ converts $B'$-coordinates to $B$-coordinates: apply $S$, then $\MM_B(T)$, then convert back with $S^{-1}$. Every invertible $S$ arises this way, since its columns form a basis.}
\end{enumerate}

\section{This week}
\subsection{Trace (Ch.~16)}

\begin{qbox}
  The matrix of an operator depends on the basis. Which numbers computed from it do not? The trace (sum of the diagonal) is one, and it equals the sum of the eigenvalues. It gives information about eigenvalues without computing them.
\end{qbox}

\begin{definition*}[Def.~41, 42 / LADR 8.47, 8.50, 8.51]
  $\tr A=a_{11}+\dots+a_{nn}$. For $T\in\LL(V)$, $\tr T=\tr\MM(T)$ in any basis.
\end{definition*}

\begin{proposition*}[Prop.~78 / LADR 8.49]
  For $A$ ($m\times n$) and $B$ ($n\times m$), $\tr(AB)=\tr(BA)$. Hence $\tr(S^{-1}AS)=\tr A$, and $\tr T$ is well defined.
\end{proposition*}
\begin{ideabox}
To show: $\tr(AB)=\tr(BA)$, although $AB\ne BA$ in general.

Write out the diagonal: $(AB)_{ii}=\sum_ja_{ij}b_{ji}$. So $\tr(AB)$ is the sum of all products $a_{ij}b_{ji}$, and $\tr(BA)$ is the sum of the same products in a different order.

Similarity: $S^{-1}AS$ is the product of $S^{-1}A$ and $S$; swap the two factors.
\end{ideabox}
\begin{proof}
  \begin{steps}
    \ii $\tr(AB)=\sum_i\sum_ja_{ij}b_{ji}=\sum_j\sum_ib_{ji}a_{ij}=\tr(BA)$.
    \ii $\tr\big((S^{-1}A)S\big)=\tr\big(S(S^{-1}A)\big)=\tr A$. Matrices of one operator are similar (T1).
  \end{steps}
\end{proof}

\begin{proposition*}[Prop.~79 / LADR 8.52]
  Over $\CC$, $\tr T=\sum_id_i\lambda_i$ over the distinct eigenvalues, $d_i=\dim G(\lambda_i,T)$. Also $-\tr T$ is the coefficient of $x^{n-1}$ in $\chi_T$.
\end{proposition*}
\begin{ideabox}
To show: $\tr T$ in terms of the eigenvalues.

$\tr T$ can be computed in any basis (Prop.~78). In a Jordan basis, or the triangular basis of LADR 8.37, the diagonal lists each $\lambda_i$ exactly $d_i$ times.

Use that basis. For $\chi_T$: expanding $\prod(x-\lambda_i)^{d_i}$, the coefficient of $x^{n-1}$ takes $-\lambda_i$ from one factor and $x$ from all the others, giving $-\sum d_i\lambda_i$.
\end{ideabox}
\begin{proof}
  In a Jordan basis the blocks for $\lambda_i$ have total size $d_i$. Expanding $\prod(x-\lambda_i)^{d_i}$, the $x^{n-1}$ coefficient is $-\sum d_i\lambda_i$.
\end{proof}

\begin{proposition*}[Prop.~80]
  $\tr(ST-TS)=0$. So on $V\ne\{0\}$ finite-dimensional, there are no $S,T$ with $ST-TS=I$.
\end{proposition*}
\begin{ideabox}
To show: The equation $ST-TS=I$ has no solution.

A standard way to show an equation has no solution: apply a function that always gives the two sides different values.

Use the trace. It is linear and does not see the order of a product, so $\tr(ST-TS)=0$ for all $S,T$, while $\tr I=\dim V\ne0$.
\end{ideabox}
\begin{proof}
  $\tr(ST-TS)=\tr(ST)-\tr(TS)=0$ (linearity, Prop.~78), while $\tr I=\dim V>0$.
\end{proof}

\begin{itemize}
  \ii \textbf{Linear:} $\tr(A+B)=\tr A+\tr B$, $\tr(cA)=c\tr A$ (diagonals add).
  \ii \textbf{Cyclic:} $\tr(ABC)=\tr(BCA)=\tr(CAB)$ (Prop.~78 with $A$ and $BC$). Not every order: $A=E_{12},B=E_{21},C=E_{11}$ give $\tr(ABC)=1$, $\tr(BAC)=0$.
  \ii \textbf{Projections:} $P^2=P\Rightarrow\tr P=\operatorname{rank}P$. In a basis for $\nul P\oplus\range P$ (Prop.~91), $\MM(P)=\operatorname{diag}(0,\dots,0,1,\dots,1)$.
  \ii \textbf{Nilpotent:} $\tr N=0$, and $\tr N^k=0$. The converse fails: $\operatorname{diag}(1,-1)$.
  \ii \textbf{Polynomials:} over $\CC$, $\tr p(T)=\sum_id_ip(\lambda_i)$ ($p(T)$ is triangular in a Jordan basis with diagonal $p(\lambda_i)$).
  \ii \textbf{$2\times2$:} $\lambda_1+\lambda_2=\tr A$, $\lambda_1\lambda_2=\det A$. Ex.~64: $\tr\begin{psmallmatrix}1&1&1\\1&2&4\\1&3&9\end{psmallmatrix}=12$ is the sum of eigenvalues that are hard to find.
  \ii \textbf{Detecting $0$:} if $\tr(ST)=0$ for all $S$, then $T=0$ (take $S=E_{ji}$: $\tr(E_{ji}\MM(T))=t_{ij}$).
\end{itemize}

\begin{trapbox}
  \ii $\tr(AB)\ne\tr A\cdot\tr B$ ($A=B=I_2$: $2\ne4$).
  \ii $\tr T=0$ does not make $T$ nilpotent.
  \ii $ST-TS=I$ is possible in infinite dimensions: on $\PP(\RR)$, $D(xp)-xD(p)=p$.
\end{trapbox}



\subsection{Inner products and norms (Ch.~17)}

\begin{qbox}
  So far vector spaces have no lengths or angles. In $\RR^n$ the dot product gives both. We want them in any vector space (polynomials, functions), so we take the properties of the dot product that proofs use as the definition.
\end{qbox}

$\FF$ is $\RR$ or $\CC$ (over $\RR$ ignore the bars). From here on $V$ may be infinite-dimensional; finite dimension is stated when needed.

\begin{definition*}[Def.~43 / LADR 6.2]
  An \vocab{inner product} $\ip--\colon V\times V\to\FF$ satisfies
  \begin{enumerate}[(i)]
    \ii positivity: $\ip vv>0$ for $v\ne0$;
    \ii linearity in the first slot: $\ip{\lambda u+v}w=\lambda\ip uw+\ip vw$;
    \ii conjugate symmetry: $\ip uv=\overline{\ip vu}$.
  \end{enumerate}
\end{definition*}
\begin{itemize}
  \ii \textbf{Why the conjugate:} without it, $(i,1)\cdot(i,1)=i^2+1=0$, a nonzero vector of length $0$. With $z\bar z=|z|^2\ge0$ lengths are positive, and symmetry must become conjugate symmetry.
  \ii $\ip0v=0$: $\ip0v=\ip{0+0}v=2\ip0v$.
  \ii Second slot: $\ip u{\lambda v+w}=\overline{\ip{\lambda v+w}u}=\bar\lambda\ip uv+\ip uw$ (conjugate-linear).
  \ii $\ip vv$ is real (it equals its own conjugate).
  \ii \textbf{Examples:} dot product $\sum w_j\bar z_j$ on $\FF^n$ (Ex.~65); $L^2$ product $\int_a^bf\bar g$ on continuous functions (Ex.~66); weighted $x_1y_1+2x_2y_2$.
  \ii \textbf{Checking a candidate:} (ii) and (iii) are usually immediate. Test positivity: can $\ip vv\le0$ for some $v\ne0$? For integrals: $\int p^2=0$ forces the continuous $p^2\ge0$ to vanish on the interval, so $p$ has infinitely many roots and $p=0$.
\end{itemize}

\begin{definition*}[Def.~44 / Prop.~81 / LADR 6.7, 6.17]
  A \vocab{norm}: $\|v\|>0$ for $v\ne0$; $\|\lambda v\|=|\lambda|\|v\|$; $\|u+v\|\le\|u\|+\|v\|$. An inner product gives $\|v\|=\sqrt{\ip vv}$.
\end{definition*}
\begin{itemize}
  \ii Homogeneity: $\|\lambda v\|^2=\lambda\bar\lambda\ip vv=|\lambda|^2\|v\|^2$. The triangle inequality needs Cauchy--Schwarz.
  \ii \textbf{Pythagoras:} if $\ip uv=0$, then $\|u+v\|^2=\|u\|^2+\|v\|^2$ (expand; the cross terms are $0$).
  \ii The converse holds over $\RR$ but not over $\CC$ (the cross term is $2\operatorname{Re}\ip uv$): in $\CC$, $u=1$, $v=i$ give $|1+i|^2=1+1$ but $\ip uv=-i$.
\end{itemize}

\begin{proposition*}[LADR 6.13: orthogonal decomposition]
  If $v\ne0$: $u=cv+w$ with $w\perp v$, where $c=\dfrac{\ip uv}{\|v\|^2}$. So the projection onto a line is $P_{\spn(v)}u=\dfrac{\ip uv}{\|v\|^2}v$.
\end{proposition*}
\begin{ideabox}
To show: Write $u$ as a multiple of $v$ plus a vector perpendicular to $v$. In $\RR^2$ this is dropping a perpendicular from $u$ to the line through $v$.

Write $u=cv+(u-cv)$; only $c$ is unknown, and it must make $u-cv\perp v$. Solve $0=\ip{u-cv}v=\ip uv-c\|v\|^2$.

Why it matters: Cauchy--Schwarz, Gram--Schmidt, orthogonal projections and best approximation all use this computation.
\end{ideabox}
\begin{fig}
\begin{tikzpicture}[scale=1.15,>=Stealth]
  \draw[gray] (-0.4,0) -- (4.2,0);
  \draw[->,thick] (0,0) -- (3.4,0) node[below right] {$v$};
  \draw[->,very thick,red] (0,0) -- (2,0);
  \node[red,below] at (1,0) {$cv$};
  \draw[->,thick,blue] (0,0) -- (2,1.6) node[above left] {$u$};
  \draw[->,thick,OliveGreen] (2,0) -- (2,1.6);
  \node[OliveGreen!70!black,right] at (2,0.8) {$w=u-cv$};
  \draw (1.82,0) -- (1.82,0.18) -- (2,0.18);
\end{tikzpicture}
\figcap{$c=\ip uv/\|v\|^2$ is the one value that makes $w\perp v$. Pythagoras gives $\|u\|^2=\|cv\|^2+\|w\|^2\ge\|cv\|^2=|\ip uv|^2/\|v\|^2$, which is Cauchy--Schwarz. Gram--Schmidt subtracts such a $cv$ for each earlier $e_k$.}
\end{fig}

\begin{theorem*}[Prop.~82 / LADR 6.14: Cauchy--Schwarz]
  $|\ip uv|\le\|u\|\|v\|$, with equality iff one of $u,v$ is a multiple of the other.
\end{theorem*}
\begin{ideabox}
To show: $|\ip uv|\le\|u\|\|v\|$.

In $\RR^2$, the projection of $u$ onto the line through $v$ is no longer than $u$, and its length is $|\ip uv|/\|v\|$. Pythagoras is the general tool for ``a perpendicular piece of a vector is no longer than the vector''.

Decompose $u=cv+w$ with $w\perp v$ (LADR 6.13). Pythagoras: $\|u\|^2=|c|^2\|v\|^2+\|w\|^2\ge|c|^2\|v\|^2=|\ip uv|^2/\|v\|^2$. Equality iff $w=0$, i.e.\ $u$ is a multiple of $v$.

Notes' version: The notes expand $\ip zz\ge0$ with $z=-w$; it is the same computation.
\end{ideabox}
\begin{proof}
  \begin{steps}
    \ii If $v=0$, both sides are $0$ and $v=0\cdot u$.
    \ii If $v\ne0$, write $u=cv+w$, $c=\ip uv/\|v\|^2$, $w\perp v$. Pythagoras: $\|u\|^2=|c|^2\|v\|^2+\|w\|^2\ge\dfrac{|\ip uv|^2}{\|v\|^2}$.
    \ii Equality iff $w=0$, i.e.\ $u=cv$. Conversely $u=av$ gives $|\ip uv|=|a|\|v\|^2=\|u\|\|v\|$.
  \end{steps}
\end{proof}

\begin{proposition*}[Prop.~81(iii) / LADR 6.17: triangle inequality]
  $\|u+v\|\le\|u\|+\|v\|$, with equality iff one is a nonnegative multiple of the other.
\end{proposition*}
\begin{ideabox}
To show: $\|u+v\|\le\|u\|+\|v\|$. Both sides are $\ge0$, so it is enough to compare squares.

$\|u+v\|^2=\|u\|^2+2\operatorname{Re}\ip uv+\|v\|^2$ and $(\|u\|+\|v\|)^2=\|u\|^2+2\|u\|\|v\|+\|v\|^2$. Only the middle terms differ.

$\operatorname{Re}\ip uv\le|\ip uv|\le\|u\|\|v\|$ by Cauchy--Schwarz. Equality needs both of these to be equalities.
\end{ideabox}
\begin{proof}
  \begin{steps}
    \ii $\|u+v\|^2=\|u\|^2+2\operatorname{Re}\ip uv+\|v\|^2\le\|u\|^2+2|\ip uv|+\|v\|^2\le(\|u\|+\|v\|)^2$.
    \ii Equality needs $\operatorname{Re}\ip uv=|\ip uv|=\|u\|\|v\|$: then $u=av$ (or $v=0$), and $\ip uv=a\|v\|^2\ge0$ forces $a\ge0$.
    \ii Conversely, $u=av$ with $a\ge0$ gives $\|u+v\|=(1+a)\|v\|=\|u\|+\|v\|$.
  \end{steps}
\end{proof}

\begin{proposition*}[Prop.~83 / LADR 6.21, 6A.28: parallelogram law, polarization]
  A norm comes from an inner product iff $\|u+v\|^2+\|u-v\|^2=2\|u\|^2+2\|v\|^2$. Then
  \[ \ip uv=\tfrac14\big(\|u+v\|^2-\|u-v\|^2\big)\ (\RR),\qquad\ip uv=\tfrac14\big(\|u+v\|^2-\|u-v\|^2+i\|u+iv\|^2-i\|u-iv\|^2\big)\ (\CC). \]
\end{proposition*}
\begin{ideabox}
To show: (a) A relation among norms that every norm from an inner product satisfies. (b) A formula that recovers $\ip uv$ from norms.

Expand $\|u+v\|^2$ and $\|u-v\|^2$: the cross term $2\operatorname{Re}\ip uv$ appears with opposite signs. Adding cancels it (the law); subtracting isolates it (polarization).

Over $\CC$: Subtracting gives only $\operatorname{Re}\ip uv$. Since $\operatorname{Re}\ip u{iv}=\operatorname{Im}\ip uv$, do the same with $iv$ in place of $v$ to get the imaginary part.

Note: The converse (the law forces an inner product) is not proved in the notes.
\end{ideabox}
\begin{proof}
  \begin{steps}
    \ii $\|u\pm v\|^2=\|u\|^2\pm2\operatorname{Re}\ip uv+\|v\|^2$. Add: the law. Subtract: $4\operatorname{Re}\ip uv$.
    \ii $\operatorname{Re}\ip u{iv}=\operatorname{Re}(-i\ip uv)=\operatorname{Im}\ip uv$, so $\|u+iv\|^2-\|u-iv\|^2=4\operatorname{Im}\ip uv$.
  \end{steps}
\end{proof}
Use: the max norm $\max(|x|,|y|)$ on $\RR^2$ fails the law ($u=e_1,v=e_2$: $1+1\ne2+2$), so it comes from no inner product.

\begin{itemize}
  \ii \textbf{Angle} (Def.~45 / LADR 6A.16, real spaces): for $u,v\ne0$, $\cos\theta=\frac{\ip uv}{\|u\|\|v\|}$, $\theta\in[0,\pi]$. Cauchy--Schwarz puts the quotient in $[-1,1]$. Ex.~68: $(1,0),(1,1)$ make angle $\pi/4$.
  \ii \textbf{Using Cauchy--Schwarz:} write the quantity as $\ip uv$ with a clever second vector, usually all $1$s or the constant $1$.
  \begin{itemize}
    \ii $(a_1+\dots+a_n)^2\le n\sum a_i^2$, equality iff all $a_i$ equal.
    \ii $\big(\int_0^1f\big)^2\le\int_0^1f^2$; $\int_0^\pi\sqrt{\sin x}\,dx\le\sqrt{2\pi}$ (Ex.~67).
    \ii Weighted: $\big(\sum a_kb_k\big)^2\le\big(\sum ka_k^2\big)\big(\sum b_k^2/k\big)$, writing $a_kb_k=(\sqrt k\,a_k)(b_k/\sqrt k)$.
  \end{itemize}
  \ii \textbf{Reverse triangle inequality:} $\big|\|u\|-\|v\|\big|\le\|u-v\|$, since $\|u\|\le\|u-v\|+\|v\|$ and symmetrically.
\end{itemize}

\needspace{7\baselineskip}\textbf{Attention (notes Section 17.3, optional application).} Stored items have keys $k_j\in\RR^{d_k}$ and values $\nu_j\in\RR^{d_v}$; a query $q\in\RR^{d_k}$ is compared with the keys.
\begin{itemize}
  \ii Scores $s_j=\ip q{k_j}/\sqrt{d_k}$; weights by softmax, $a_j=e^{s_j}/\sum_\ell e^{s_\ell}$ (so $a_j>0$, $\sum a_j=1$); output $y=\sum_ja_j\nu_j$. Keys decide how much each item contributes; values are what is returned. $d_k$ and $d_v$ can differ.
  \ii Ex.~69: $k_1=(2,0)$, $k_2=(0,1)$, $k_3=(-1,0)$, values $e_1,e_2,e_3\in\RR^3$, $q=(1,0)$, $d_k=2$. Scores $\sqrt2,0,-\frac1{\sqrt2}$; weights $\approx0.734,0.178,0.088$; $y\approx0.734e_1+0.178e_2+0.088e_3$.
  \ii $\ip qk=\|q\|\|k\|\cos\theta$, so the score depends on lengths as well as the angle. Why $\sqrt{d_k}$: if the coordinates $q_r,k_r$ are independent with mean $0$ and variance $1$, then $\ip qk=\sum_{r=1}^{d_k}q_rk_r$ has variance $d_k$; dividing by $\sqrt{d_k}$ makes the variance $1$.
  \ii Self-attention: $q_i=W_Qx_i$, $k_j=W_Kx_j$, $\nu_j=W_Vx_j$ for learned linear maps. With $Q,K,V_{\mathrm{val}}$ having columns $q_i,k_j,\nu_j$: score matrix $S=Q^{\mathsf T}K/\sqrt{d_k}$, softmax of each row gives $A=(a_{ij})$, which is row-stochastic, and $Y=V_{\mathrm{val}}A^{\mathsf T}$ has columns $y_i=\sum_ja_{ij}\nu_j$. For left-to-right generation, terms with $j>i$ are removed before softmax (masking). Because $A$ depends on the inputs and softmax is nonlinear, self-attention is not a linear map.
\end{itemize}

\begin{trapbox}
  \ii Over $\CC$, scalars leave the second slot conjugated.
  \ii $x_1y_1-x_2y_2$ and $(x_1+x_2)(y_1+y_2)$ are not inner products: $\ip{e_2}{e_2}=-1$, $\ip{(1,-1)}{(1,-1)}=0$.
  \ii Cauchy--Schwarz equality: ``dependent''. Triangle equality: ``nonnegative multiple''.
\end{trapbox}



\section{Problem set}

\begin{prob}
\textbf{HW9 P1. }\ Let $P$ be a projection on a finite-dimensional vector space, so that $P^2=P$. Prove that $\tr P=\operatorname{rank}P$. You may choose a basis adapted to the decomposition $V=\nul P\oplus\range P$.
\end{prob}

\begin{prob}
\textbf{HW9 P2. }\ Let $T\in\LL(\CC^5)$ have minimal polynomial $m_T(x)=(x-2)^2(x+1)$.
\begin{enumerate}[(a)]
  \ii Determine all possible algebraic multiplicities of the eigenvalues $2$ and $-1$.
  \ii Determine all possible values of $\tr T$.
  \ii For each possible trace, exhibit a Jordan normal form realizing it.
\end{enumerate}
\end{prob}

\begin{prob}
\textbf{HW9 P3. }\ 
\begin{enumerate}[(a)]
  \ii Prove that every nilpotent operator on a finite-dimensional complex vector space has trace $0$.
  \ii Give a nonzero operator with trace $0$ that is not nilpotent.
  \ii Give two $2\times2$ matrices $A,B$ for which $AB-BA\ne0$, and verify directly that $\tr(AB-BA)=0$.
\end{enumerate}
\end{prob}

\begin{prob}
\textbf{Tr1. }\ Let $A$ be a real $2\times2$ matrix with $A^2=-I$. Prove that $\tr A=0$ and $\det A=1$.
\end{prob}

\begin{prob}
\textbf{Tr2. }\ Let $T\in\LL(\CC^3)$ have $m_T=(x-1)(x-2)^2$. Find $\chi_T$, the Jordan form of $T$, and $\tr T$.
\end{prob}

\begin{prob}
\textbf{HW9 P4. }\ For each of the following formulas on $\RR^2$, decide whether it defines an inner product. Justify your answer.
\[ \ip xy_1=x_1y_1+2x_2y_2,\qquad\ip xy_2=x_1y_1-x_2y_2,\qquad\ip xy_3=(x_1+x_2)(y_1+y_2). \]
\end{prob}

\begin{prob}
\textbf{HW9 P5. }\ On $\PP_2(\RR)$ define $\ip pq=\int_{-1}^1p(x)q(x)\,dx$.
\begin{enumerate}[(a)]
  \ii Prove that this is an inner product.
  \ii Compute the norms of $1$, $x$, and $x^2$.
  \ii Determine which pairs among $1,x,x^2$ are orthogonal.
\end{enumerate}
\end{prob}

\begin{prob}
\textbf{HW9 P6. }\ Use the Cauchy--Schwarz inequality to prove that for all real numbers $a,b,c$, $(a+b+c)^2\le3(a^2+b^2+c^2)$. Determine exactly when equality holds.
\end{prob}

\begin{prob}
\textbf{HW9 P7. (optional)}\ Prove the parallelogram identity $\|u+v\|^2+\|u-v\|^2=2\|u\|^2+2\|v\|^2$ in every real or complex inner-product space.
\end{prob}

\begin{prob}
\textbf{I1. }\ Let $a_1,\dots,a_n>0$. Prove that $(a_1+\dots+a_n)\big(\frac1{a_1}+\dots+\frac1{a_n}\big)\ge n^2$, and find when equality holds.
\end{prob}

\end{document}
